% Nastavení dokumentu
\documentclass[a4paper,12pt]{report}
\usepackage[utf8]{inputenc}
\usepackage{hyperref}
\usepackage{graphicx}
\usepackage[table]{xcolor}
\usepackage{tabularray}
\usepackage{geometry}

% Vlastní příkazy
\newcommand{\chp}[1]{\chapter*{#1}\addcontentsline{toc}{chapter}{#1}}
\newcommand{\schp}[1]{\section*{#1}\addcontentsline{toc}{section}{#1}}

% Nastavení Titulní strana: Nadpis, autor, datum
\title{Laboratorní cvičení ADC}
\author{Vojtěch Vašek}
\date{\today}

\begin{document}

% Tisk titulní strany, obsahu a vložení nové stránky
\maketitle
\newgeometry{top=0cm}\tableofcontents\restoregeometry
\newpage

\newgeometry{top=0cm}
\chp{Zadání a úvaha}

Cílem je změna vzorkovacího kanálu.

V pomocném zdroji je AD převodník nastaven na kanál free running. Free running konfigurace znamená, že jakmile se dokončí konverze vzorku a jeho uložení do \verb|ADCH:ADCL| registru, okamžitě začíná konverze dalšího vzorku.

Další typy konfigurace převodníku jsou single conversion a diferenciální převod.

Single conversion je když po vytvoření jednoho vzorku periferie počká na pokyn k započetí další konverze.

Difereciální konfigurace umožní vzorkování dvou vstupních signálů tak, že výsledný vzorek je rozdíl mezi hladinami těchto dvou signálů v daném časovém kvantu.

\chp{Použité přístroje}
\begin{enumerate}
	\item[•]{Programátor - propojení desky s počítačem; AVR-ISP}
	\item[•]{deska s ATmega16 - typ krystalu: HC49/S QM - 14.7456 MHz, pin A7, potenciometr}
	\item[•]{Počítač - Windows 10 21H2; Kate; AVR Studio 4; AVR\_ISP\_prog}
\end{enumerate}
\restoregeometry
\schp{Obrázek zapojení}
\begin{center}
\includegraphics[width=11cm]{MIT\_14\_3\_25.jpg}
\end{center}

\newgeometry{top=0cm}
\chp{Postup práce}

V Kate jsem si napsla kód programu, který při projetím \verb|Main| smyčky zvýší hodnotu čítače (Timer\_Counter\_0). Při každém zvýšení hodnoty v registru čítače se hodnota porovná hodnota čítače (TCNT0) s maximální TOP (OCR0) hodnotou. Pokud dojde k přetečení (TCNT0 $>$ OCR0), čítač pošle jádru žádost o přerušení (TIM0\_COM).

V obsluze přerušení se mi zvíší hodnota počtu přerušení o jedno, zavolá se \verb|Toggle| podprogram, který mi na Portu A přepne hodnotu nultého bitu a vynuluje TCNT.

Celý tento kód jsem zkompiloval a ozkoušel jsem v AVR Studiu.

Následně při měření jsem v AVR Studiu měnil TOP hodnotu a hodnotu dělení hodin, upravený program jsem sestavil, přes AVR\_ISP\_prog a programátor jsem nahrál na desku a z PCLab2000LT jsem četl frekvenci osciloskopu.

Frekvenci jsem si zapsal do tabulky, která mi automaticky spočítala chybu měření oproti vypočtené frekvenci.

\schp{Zdrojový kód}
Pro lepší přehlednost tohoto dokumentu jsem se rozhodl nevkládat zdrojový kód přímo zde, ale odkázat na něj pomocí \href{https://github.com/Honeymoon-with-Anxiety/Honeymoon/tree/E4A/MIT/6_12_24}{URL}.

Kód jsem si rozdělil na \verb|templateM16|, kde mám \verb|RESET| obsluhu, hlavičku \verb|display|, kde mám podprogramy zařizující funkčnost displaye a nakonec \verb|6_12_24|, kde mám \verb|Main| smyčku.

\newgeometry{top=0cm}
\chp{Závěr}
Po otočením potenciometrem se mi na display ukáže číselná hodnota v desítkové soustavě.
\schp{Co jsem se naučil}
Jak nakonfigurovat ADC pro ATmega16.
\schp{Co jsem doufal, že se naučím a nenaučil}
Od této úlohy jsem nevěděl, co očekávat a tak jsem nemohl tušit, co se všechno naučím.
\schp{Co jsem nestihl}
Přijít na to, proč se mi na display neukazují čísla od nuly.
\schp{Jak bych pokračoval, kdybych měl více času}
Zkusil bych změnit konfiguraci převodníku v \verb|RESET| obsluze, dále bych pomocí měřících přístrojů změřil napětí v různých hodnotách a hledal mezi nimi spojitost. V nejhorším případě bych požádal o pomoc.
\restoregeometry

\end{document}
